\documentclass[a4paper,11pt]{article}

\usepackage{revy}
\usepackage[utf8]{inputenc}
\usepackage[T1]{fontenc}
\usepackage[danish]{babel}


\revyname{MatematikRevyen}
\revyyear{2026}
% HUSK AT OPDATERE VERSIONSNUMMER
\version{0.1}
\eta{$3$ minutter}
\status{Færdig}

\title{Funktionernes Pride}
\author{Oliver 25', Ronja 25', Flemming 25'}

\begin{document}
\maketitle

\begin{roles}
\role{K}[Kommentator] Announcerer der gennemgår paraden og fortæller jokes/puns
\end{roles}

\begin{props}
\prop{Rekvisit}[Person, der skaffer] En hel masse tegninger af afbildninger i diverse farver, ret rekvisit tungt
\end{props}


\begin{sketch}

\scene{Scenen starter tomt, og så træder \textbf{K} ind på scenen}

\says{K} Store UP1, lad mig høre noget jubel for Funktionernes HBTQ Pride Parade!! 

\says{K} Vi er i dag samlet for at fejre alle de smukke afbildninger vi arbejder med hver dag her på matematik, og alle de forskellige former som de findes i. Jeg kan slet ikke vente - skal vi ikke bare gå i gang? \act{træder til scene venstre}

\says{K} Vi starter naturligvis med vores "H" afbildninger, nemlig homomorfierne!

\scene{Statister træder ind på scenen med illustrationer af homomorfier i regnbue farver}

\says{K} Kendetegnet ved en homomorfi er jo at de er til deres egen gruppestruktur, til trods for så mange andre afbildninger. Men det som der ofte overses, det som der egentlig er kernen ved en homomorfi, er alle de elementer, der allesammen skaber identiteten.

\scene{Homomorfi statister bevæger sig af scenen}

\says{K} Next up, giv mig et "B"  for bijektion!

\scene{Statister træder ind med eksempler på bijektive funktioner i farverne fra bisexual pride flaget}

\says{K} Bijektioner er jo nogle af de allerfrækkeste funktioner der findes \act{catcall fløjt}. De er nemlig både til det injektive og det surjektive - de kan både lide at fylde, og selv blive fyldt ud, if you catch my drift.

\scene{Statisterne træder af scenen}

\says{K} Nu kommer vore dejlige dejlige "T" funktioner - translationer! 

\scene{Statister træder ind med translation illustrationer i trans pride flag farverne. Statisterne skal kun bevæge i lige linjer fra punkt til punkt}

\says{K} Det er faktisk sådan, at translationerne startede et heeeelt andet sted, nogen af dem stammer faktisk helt fra origo! Men nu hvor de har fået tid til at udfolde sig, så er de endt et sted hvor de føler sig meget mere hjemme.

\scene{Statisterne træder af scenen, igen kun i lige bevægelser}

\says{K} Nå, så er vi ved at være nået til slutningen af vores parade, og jeg siger tusind tak til alle jer for at have været så godt et publikum. Vi mangler jo selvfølgelig vores aller sidste bogstav "Q", og som afslutning kunne det jo ikke være andet: - \act{\textbf{K} hiver en hvid firkant ud af lommen} - QED.

\scene{Tæppe}

\end{sketch}
\end{document}
